%%==============================================================================
%% Varios formatos | @file /informe_latex/preamble.tex
%%==============================================================================
%% Lo que podria querer configurar:
%% -- graphicspath
%%==============================================================================


% document parameters
\usepackage[spanish, es-tabla]{babel}

% Uso margenes top y bottom para ajuste si quiero una cantidad de hojas.
\usepackage[vmargin = 1.9cm, hmargin = 2.1cm]{geometry}


% fgr
\usepackage[utf8]{inputenc}
\usepackage[T1]{fontenc}
\usepackage{hyperref}

% Packages for math
\usepackage{mathrsfs}
\usepackage{amsfonts}
\usepackage{amsmath}
\usepackage{amsthm}
\usepackage{amssymb}
\usepackage{physics}
\usepackage{esint}

% Packages for writing
\usepackage{enumerate}
\usepackage[shortlabels]{enumitem}
\usepackage{framed}
\usepackage{csquotes}

% Miscellaneous packages
\usepackage{float}
\usepackage{tabularx}
\usepackage{xcolor}

\usepackage{xcolor}
\usepackage{colortbl} % <-- AGREGA ESTA LÍNEA DEBAJO DE XCOLOR

\usepackage{multicol}
\usepackage{subcaption}
\usepackage{caption}

\usepackage[output-decimal-marker={,},
per-mode=symbol-or-fraction,
retain-unity-mantissa=true,
detect-all,
range-phrase={ -- },
list-final-separator={ y }
]{siunitx}

% Para imagenes usar path por defecto:
% <latex_fig src="02_perfil.png" cap="Error regresión lineal." lbl="fig:error"/>
\usepackage{graphicx}
\graphicspath{{../img/}}
\captionsetup{format = hang, margin = 10pt, font = small, labelfont = bf}

\usepackage{titlesec}
\usepackage{listings}
\usepackage[many]{tcolorbox}

% Citation
\usepackage[round, authoryear]{natbib}

% Hyperlinks setup
\usepackage{hyperref}
\definecolor{links}{rgb}{0.36,0.54,0.66}
\hypersetup{
	colorlinks = true,
	linkcolor = black,
	urlcolor = blue,
	citecolor = blue,
	filecolor = blue,
	pdfauthor = {Author},
	pdftitle = {Title},
	pdfsubject = {subject},
	pdfkeywords = {one, two},
	pdfproducer = {LaTeX},
	pdfcreator = {pdfLaTeX},
}

% Colores para las tablas.
\definecolor{tableHeaderColor}{rgb}{0.88, 0.88, 0.88}
\definecolor{tableColor}{rgb}{0.93, 0.93, 0.93}

% Para codigo.
\definecolor{pyKeyword}{rgb}{0.0, 0.5, 0.0}     % Verde para def, import, in
\definecolor{pyLib}{rgb}{0.75, 0.45, 0.0}       % Marrón/Naranja para tus variables y np
\definecolor{pyFunc}{rgb}{0.0, 0.35, 0.75}      % Azul para funciones analíticas
\definecolor{pyString}{rgb}{0.7, 0.1, 0.1}      % Rojo para textos
\definecolor{pyComment}{rgb}{0.4, 0.4, 0.4}     % Gris para comentarios
\definecolor{pyBrackets}{rgb}{0.75, 0.3, 0.65}
\definecolor{pyNumbers}{rgb}{0.3, 0.75, 0.3}     % Verde claro para números (Estilo VS Code)

\definecolor{codeBlockColor}{rgb}{0.93, 0.93, 0.93}

\lstset{
language=Python,               
basicstyle=\ttfamily\small,
columns=fixed,                 
basewidth=0.475em,             
breaklines=true,
backgroundcolor=\color{codeBlockColor}, 
frame=single,
rulecolor=\color[rgb]{0, 0, 0}, 
postbreak=\mbox{\textcolor{gray}{$\hookrightarrow$}\space},
showstringspaces=false,
% 1. Sintaxis base de Python
keywordstyle=\color{pyKeyword}\bfseries,
commentstyle=\color{pyComment}\itshape,
stringstyle=\color{pyString},
% -------------------------------------------------------------
% FORCE CON ENFASIS: Colorea funciones analíticas sin importar los puntos
emph={polyfit, log, copy, print, len, range, plot, where, diff, argmin, max, min, mean, sqrt, exp, linspace, savefig, xlabel, ylabel, xlim, ylim, grid, legend, show},
emphstyle=\color{pyFunc}\bfseries,
% -------------------------------------------------------------
% FORCE CON ENFASIS: Colorea alias de librerías y dataframes principales
emph={[2]np, pd, plt},
emphstyle={[2]\color{pyLib}\bfseries},
% -------------------------------------------------------------
% LITERATE: Mapeo exhaustivo para forzar la colorización de paréntesis cerrados
literate=*{[}{{\textcolor{pyBrackets}{[}}}1
			{]}{{\textcolor{pyBrackets}{]}}}1
			{\{}{{\textcolor{pyBrackets}{\{}}}1
			{\}}{{\textcolor{pyBrackets}{\}}}}1
}

% Silenciar la advertencia de conflicto entre siunitx y physics
\ExplSyntaxOn
\msg_redirect_name:nnn {siunitx} {physics-pkg} {none}
\ExplSyntaxOff

% fgr
\captionsetup{hypcap=false}
\usepackage[utf8]{inputenc}
\usepackage[T1]{fontenc}

% Ajuste del margen de maniobra tipográfico
\emergencystretch=3em  % Espacio extra de emergencia si una línea no calza bien
%\tolerance=2000        % Permite mayor estiramiento antes de dar una advertencia
%\hfuzz=2pt             % Ignora advertencias si el desborde es menor a 2pt (~0.7 mm)
\hbadness=10000        % Silencia TODAS las advertencias de Underfull \hbox (espaciado estirado)
%\vbadness=10000        % Silencia las advertencias de Underfull \vbox (espacio vertical suelto)

\usepackage{pgffor} % Asegurate de tener esto en tu preámbulo

% REDEFINICIÓN DE CAPTION PARA FIGURAS
\addto\captionsspanish{\renewcommand{\figurename}{Imagen}}

\usepackage{microtype} % Para angostar fuente de secciones de codigo.

% =========================================================================
% BLOQUE DE PRUEBA DE COLORES - SOLUCIÓN DEFINITIVA SINCRO
% =========================================================================
\usepackage{etoolbox}

% Para tablas, padding
% 1. Global Horizontal Padding control (Default is usually 6pt)
\setlength{\tabcolsep}{4pt} 
% 2. Fix for Nested Headers padding buildup
\usepackage{array}
\newcommand{\headercell}[1]{\begin{tabular}{@{}c@{}}#1\end{tabular}}

%%%%%%%%%%%%%%%%%%%%%%%%%%%%%%%%%%%%%%%%%%%%
% =============================================================================
% Configuración Global de Notas al Pie con Formato [1]
% =============================================================================
\makeatletter
% Ponemos corchetes al número de llamada en el texto principal
\renewcommand\@makefnmark{\hbox{\@textsuperscript{\normalfont[\@thefnmark]}}}

% Ponemos corchetes y alineamos el número abajo en el pie de página
\renewcommand\@makefntext[1]{%
    \parindent 1em%
    \noindent
    \hb@xt@1.8em{\hss\normalfont[\@thefnmark]}#1}
\makeatother

\usepackage{mathtools} % Activa \xRightarrow de forma nativa
